\documentclass{article}
\usepackage[utf8]{inputenc}
\usepackage{amsmath,amssymb}
\usepackage[most]{tcolorbox}
\usepackage{geometry}
\geometry{margin=2.5cm}

\newcommand{\step}[1]{\par\noindent\hangindent=3.5em\hangafter=1%
  \makebox[3.5em][l]{\textbf{#1.}}}
\newcommand{\steptag}[1]{\hfill{\small\textsf{[#1]}}}

\newtcolorbox{proofbox}[1]{
  colback=gray!5, colframe=gray!60,
  fonttitle=\bfseries, title={#1},
  breakable, enhanced,
  left=4pt, right=4pt, top=4pt, bottom=4pt,
  before skip=10pt, after skip=10pt
}

\begin{document}

\begin{proofbox}{There is a universal e\_zero > 0 such that, if A is a fini...}
\step{1} There is a universal e\_zero > 0 such that, if A is a finite-dimensional extended epsilon\_A-C*-algebra with epsilon\_A <= t, a supplied MAIN partition state has w:C\textasciicircum{}m->A a non-unital extended t-inclusion with one-dimensional images P\_j, U,V are disjoint nonempty unions sharing no equivalence class, R = U union V, and t <= e\_zero, then dim S\textasciicircum{}A\_\{P\_U,P\_V\} = dim S\textasciicircum{}A\_\{P\_V,P\_U\} = 0 and dim S\textasciicircum{}\{A\_R\}\_\{P\_U\textasciicircum{}R,P\_V\textasciicircum{}R\} = dim S\textasciicircum{}\{A\_R\}\_\{P\_V\textasciicircum{}R,P\_U\textasciicircum{}R\} = 0. \steptag{validated} \\[6pt]
\step{1.1} Let e\_add>0 be a universal accuracy threshold at which lem-extcb-corner-dimension-additivity applies, and set e\_zero=min\{e\_sim,e\_add,e\_ncd,1/2\}, where e\_sim and e\_ncd are the universal constants in lem-maincb-corner-equivalence and lem-maincb-nested-corner-dimension-transport. Then e\_zero>0 is universal, and t<=e\_zero implies all three external lemmas are applicable and t<1. \steptag{validated} \\[4pt]
\step{1.2} Apply lem-maincb-corner-equivalence to the one-dimensional t-projections (P\_j): P\_j=w(e\_j) is a t-projection because e\_j\textasciicircum{}2=e\_j and the extended t-inclusion has t-multiplicative defect. Hence the relation j\textasciitilde{}k iff dim S\textasciicircum{}A\_\{P\_j,P\_k\}=1 is an equivalence relation. Since U and V are disjoint unions sharing no equivalence class, for every j in U and k in V one has j not\textasciitilde{}k and k not\textasciitilde{}j, and therefore dim S\textasciicircum{}A\_\{P\_j,P\_k\} != 1 and dim S\textasciicircum{}A\_\{P\_k,P\_j\} != 1. No zero-dimensional conclusion follows from the allowed inputs alone. \steptag{validated} \\[6pt]
\step{1.2.1} For every j, P\_j=w(e\_j) is a t-projection: the projection identity e\_j\textasciicircum{}2=e\_j and the t-multiplicative-defect bound for w give ||P\_j\textasciicircum{}2-P\_j||<=t, while the self-adjointness clause for the inclusion gives the required Hermitian condition. Since t<=e\_zero<=e\_sim, lem-maincb-corner-equivalence applies and makes j\textasciitilde{}k iff dim S\textasciicircum{}A\_\{P\_j,P\_k\}=1 an equivalence relation. \steptag{validated} \\[4pt]
\step{1.2.2} If j is in U and k is in V, then j and k lie in distinct equivalence classes because U and V are disjoint unions sharing no class. Thus j not\textasciitilde{}k and, by symmetry, k not\textasciitilde{}j. By the defining biconditional for \textasciitilde{}, dim S\textasciicircum{}A\_\{P\_j,P\_k\}!=1 and dim S\textasciicircum{}A\_\{P\_k,P\_j\}!=1. This does not imply either dimension is zero: excluding dimensions at least two requires the separate one-dimensional-corner-dimension theorem, which is not an allowed external input in this workspace. \steptag{validated} \\[4pt]
\step{1.3} Restrict w to the coordinate commutative C*-algebras C\textasciicircum{}U and C\textasciicircum{}V. These restrictions remain non-unital extended t-inclusions, take their units to P\_U and P\_V, and take their projection bases to (P\_j)\_\{j in U\} and (P\_k)\_\{k in V\}. By lem-extcb-corner-dimension-additivity and the atomwise zero dimensions, dim S\textasciicircum{}A\_\{P\_U,P\_V\}=sum\_\{j in U,k in V\} dim S\textasciicircum{}A\_\{P\_j,P\_k\}=0; interchanging U,V gives dim S\textasciicircum{}A\_\{P\_V,P\_U\}=0. \steptag{validated} \\[4pt]
\step{1.4} Put P\_R=w(e\_R)=P\_U+P\_V for R=U union V. Because epsilon\_A<=t, A satisfies the extended t-C*-bounds. The elements P\_U,P\_V,P\_R are t-projections; P\_R is nonvanishing; and all four left/right subordination errors of P\_U,P\_V to P\_R are at most t. The supplied def-maincb-partition-state gives A\_R=S\textasciicircum{}A\_\{P\_R\}. Introduce locally P\_U\textasciicircum{}R:=Co\textasciicircum{}A\_\{P\_R\}(P\_U) and P\_V\textasciicircum{}R:=Co\textasciicircum{}A\_\{P\_R\}(P\_V); these are the exact compression equalities required by lem-maincb-nested-corner-dimension-transport, rather than notation asserted by def-maincb-partition-state. \steptag{validated} \\[6pt]
\step{1.4.1} Let e\_R be the coordinate projection of C\textasciicircum{}m associated to the nonempty set R=U union V, and put P\_R=w(e\_R). At level one, the extended t-inclusion w is a t-homomorphism and satisfies the two-sided (1 plus-or-minus t) norm bounds. Since e\_R is Hermitian and e\_R\textasciicircum{}2=e\_R, involution preservation and the multiplicative-defect axiom give P\_R\textasciicircum{}dagger=P\_R and ||P\_R\textasciicircum{}2-P\_R||<=t||e\_R||\textasciicircum{}2=t; hence P\_R is a t-projection by def-delta-projection. Also ||e\_R||=1, so 1-t<=||P\_R||<=1+t and therefore | ||P\_R||-1 |<=t<=t+epsilon\_A. This is the second alternative in def-delta-projection with explicit universal O-constant 1, so P\_R is nonvanishing in the sense required by lem-maincb-nested-corner-dimension-transport. \steptag{validated} \\[6pt]
\step{1.4.1.1} At amplification n=1, def-extended-delta-inclusion says that w is a t-homomorphism. Thus w preserves the involution and ||w(xy)-w(x)w(y)||<=t||x||||y||. For the coordinate projection e\_R one has e\_R\textasciicircum{}dagger=e\_R, e\_R\textasciicircum{}2=e\_R, and, because R is nonempty, ||e\_R||=1. Consequently P\_R\textasciicircum{}dagger=w(e\_R)\textasciicircum{}dagger=w(e\_R\textasciicircum{}dagger)=P\_R and ||P\_R\textasciicircum{}2-P\_R||=||w(e\_R)w(e\_R)-w(e\_R\textasciicircum{}2)||<=t. Hence P\_R is a t-projection. The two-sided norm bounds of the same inclusion give 1-t<=||P\_R||<=1+t, whence | ||P\_R||-1 |<=t<=t+epsilon\_A. Therefore P\_R satisfies the nonvanishing second alternative of def-delta-projection with universal coefficient 1, independently of the pending parent node. \steptag{validated} \\[4pt]
\step{1.4.2} Let e\_U,e\_V be the coordinate projections of C\textasciicircum{}m associated to U,V. At level one, def-extended-delta-inclusion says that w preserves the involution and has t-multiplicative defect, so P\_U=w(e\_U) and P\_V=w(e\_V) are t-projections. Since e\_Ue\_R=e\_Re\_U=e\_U and e\_Ve\_R=e\_Re\_V=e\_V, the four quantities ||P\_UP\_R-P\_U||, ||P\_RP\_U-P\_U||, ||P\_VP\_R-P\_V||, and ||P\_RP\_V-P\_V|| are at most t. The supplied def-maincb-partition-state gives A\_R=S\textasciicircum{}A\_\{P\_R\}. Introduce the local notation P\_U\textasciicircum{}R:=Co\textasciicircum{}A\_\{P\_R\}(P\_U) and P\_V\textasciicircum{}R:=Co\textasciicircum{}A\_\{P\_R\}(P\_V); these are exactly the defining equalities required as hypotheses by lem-maincb-nested-corner-dimension-transport. No independent claim that compression preserves approximate projections is made or needed: the transport lemma itself is the permitted result whose conclusion is the asserted corner-dimension equality for these compressed elements. \steptag{validated} \\[4pt]
\step{1.5} Apply lem-maincb-nested-corner-dimension-transport first to (P\_R,P\_U,P\_V) and then to (P\_R,P\_V,P\_U), using the hypotheses established above. Together with the two ambient zero dimensions, this gives dim S\textasciicircum{}\{A\_R\}\_\{P\_U\textasciicircum{}R,P\_V\textasciicircum{}R\}=0 and dim S\textasciicircum{}\{A\_R\}\_\{P\_V\textasciicircum{}R,P\_U\textasciicircum{}R\}=0. Hence all four dimension equalities asserted by node 1 hold with the universal e\_zero chosen above. \steptag{validated} \\[4pt]
\step{1.6} Shrink the final constant further to e\_zero=min\{e\_sim,e\_add,e\_ncd,e\_rank,1/2\}, where the universal e\_rank is chosen below. For any atom pair P=P\_j,Q=P\_k, the registered one-dimensional-projection definition gives dim S\_P=dim S\_Q=1. We prove dim S\_\{P,Q\}<=1 using only the registered compression estimates. All O(t+epsilon\_A) bounds below have universal constants and are O(t) since epsilon\_A<=t. If S\_\{P,Q\}=0 there is nothing to prove. Otherwise choose X in S\_\{P,Q\} with ||X||=1 and put p\textasciitilde{}=Co\_P(P), q\textasciitilde{}=Co\_Q(Q). Because P,Q are images of norm-one coordinate projections under the extended t-inclusion, ||P|| and ||Q|| are bounded above and below by 1+O(t). From Z=Co\_\{P,Q\}(Z), the estimates of Co\_\{P,Q\} against both L\_P R\_Q and R\_Q L\_P, approximate associativity, and ||P\textasciicircum{}2-P||,||Q\textasciicircum{}2-Q||<=t give, uniformly for Z in S\_\{P,Q\}, ||PZ-Z||+||ZQ-Z||<=C\_1 t||Z||. The same compression estimates give ||p\textasciitilde{}-P||+||q\textasciitilde{}-Q||<=C\_2t and hence ||p\textasciitilde{} dot Z-Z||+||Z dot q\textasciitilde{}-Z||<=C\_3t||Z||, where dot is the registered compressed product. Comparing each compressed product with the ambient product and using ambient approximate associativity also gives ||(Y dot W) dot Z-Y dot(W dot Z)||<=C\_4t||Y||||W||||Z|| for the matching corners used here. Now B=X dot X\textasciicircum{}dagger lies in the one-dimensional S\_P, so B=beta p\textasciitilde{}. The C*-lower bound applied to X\textasciicircum{}dagger gives ||XX\textasciicircum{}dagger||>=(1-epsilon\_A)||X||\textasciicircum{}2, while compressed-product closeness gives ||B-XX\textasciicircum{}dagger||<=C\_5t. Since ||p\textasciitilde{}|| is universally bounded, after decreasing e\_rank we obtain |beta|>=c>0. For arbitrary Z in S\_\{P,Q\}, write X\textasciicircum{}dagger dot Z=gamma(Z)q\textasciitilde{}; gamma is linear and |gamma(Z)|<=C\_6||Z|| because ||q\textasciitilde{}|| is bounded below. The preceding unit and associator estimates yield ||beta Z-gamma(Z)X||<=C\_7t||Z||. Thus T:Z↦beta\textasciicircum{}\{-1\}gamma(Z)X has rank at most one and ||I-T||<=C\_7t/c. Choose the universal e\_rank so that C\_7e\_rank/c<1. Then T is invertible by the Neumann lemma, and therefore dim S\_\{P,Q\}<=rank(T)<=1. The adjoint identity Co\_\{P,Q\}(Y)\textasciicircum{}dagger=Co\_\{Q,P\}(Y\textasciicircum{}dagger) is a conjugate-linear bijection S\_\{P,Q\}->S\_\{Q,P\}, so the reverse corner also has dimension at most one. Applying this to every j in U,k in V and combining it with validated node 1.2.2, which excludes dimension one, gives dim S\_\{P\_j,P\_k\}=dim S\_\{P\_k,P\_j\}=0. The allowed corner-dimension additivity lemma now gives both ambient union-corner dimensions zero. Finally, validated node 1.4 supplies every hypothesis of nested-corner dimension transport, so applying that allowed lemma in both orders transports these two zeros to A\_R and proves the four conclusions of node 1. \steptag{validated} \\[4pt]
\end{proofbox}

\end{document}
